\documentclass[11pt]{article}

% Shared notes settings for Applied Machine Learning lectures (UPES).
% Keep this file free of \documentclass to allow reuse via \input.

\usepackage[utf8]{inputenc}
\usepackage[T1]{fontenc}
\usepackage{lmodern}          % scalable T1 fonts; without this pdflatex falls back
\input{glyphtounicode}        % to bitmapped EC Type 3 and the PDFs are neither
\pdfgentounicode=1            % searchable nor screen-reader friendly
\usepackage[a4paper,margin=1in]{geometry}
\usepackage{hyperref}
\usepackage{xurl}
\usepackage{booktabs}
\usepackage{amsmath}
\usepackage{amssymb}
\usepackage{listings}
\usepackage{xcolor}
\usepackage{enumitem}
\usepackage{parskip}

% Code listings (Python-oriented for this subject).
\lstset{
  language=Python,
  basicstyle=\ttfamily\small,
  keywordstyle=\color{blue},
  commentstyle=\color{gray},
  stringstyle=\color{teal},
  showstringspaces=false,
  columns=fullflexible,
  frame=single,
  framerule=0.2pt,
  breaklines=true
}

\setlist[itemize]{topsep=0.3em,itemsep=0.2em}
\setlist[enumerate]{topsep=0.3em,itemsep=0.2em}

% Simple per-section source line.
\newcommand{\notesources}[1]{\par\noindent{\small\color{gray}\textit{Sources:} \sloppy #1\par}}

% Unit-wise source strings for this subject (used by slides + notes).
% Path is relative to the lecture `latex/` folder (the compilation CWD).
\IfFileExists{../../../_shared/aml-sources.tex}{%
  % Source strings for Applied Machine Learning (CSAI2017P) — used by slides + notes.
% Prescribed texts and reference books are taken from the course syllabus.
% Support texts are the author-hosted open-access editions:
%   statlearning.com            (ISL with Python)
%   probml.github.io/pml-book   (Murphy, Probabilistic Machine Learning)
%   d2l.ai                      (Dive into Deep Learning)
%   mml-book.github.io          (Mathematics for Machine Learning)

\newcommand{\AMLRepoURL}{https://github.com/tali7c/Applied-Machine-Learning}

% --- prescribed -------------------------------------------------------------
\newcommand{\AMLPrescribed}{%
M\"uller and Guido, \textit{Introduction to Machine Learning with Python}, O'Reilly, 2016.;%
Bishop, \textit{Pattern Recognition and Machine Learning}, Springer, 2006.;%
Murphy, \textit{Machine Learning: A Probabilistic Perspective}, MIT Press, 2012.;%
Han, Kamber and Pei, \textit{Data Mining: Concepts and Techniques} (3e), Morgan Kaufmann, 2012.%
}

% --- unit-wise ---------------------------------------------------------------
\newcommand{\AMLUnitISources}{%
M\"uller and Guido, \textit{Introduction to Machine Learning with Python}, Ch. 1.;%
Bishop, \textit{Pattern Recognition and Machine Learning}, Ch. 1.;%
Murphy, \textit{Probabilistic Machine Learning: An Introduction}, MIT Press, 2022, Ch. 1.;%
Mitchell, \textit{Machine Learning}, McGraw Hill, 1997, Ch. 1.;%
Scikit-learn documentation, accessed 2026-08-04.%
}

\newcommand{\AMLUnitIISources}{%
Bishop, \textit{Pattern Recognition and Machine Learning}, \S1.5, \S4.3, \S7.1.;%
Zhang, Lipton, Li and Smola, \textit{Dive into Deep Learning}, \S3.4.;%
Shalev-Shwartz and Ben-David, \textit{Understanding Machine Learning}, Ch. 15.;%
Murphy, \textit{Probabilistic Machine Learning: An Introduction}, Ch. 5.%
}

\newcommand{\AMLUnitIIISources}{%
Zhang, Lipton, Li and Smola, \textit{Dive into Deep Learning}, Ch. 12 (Optimization Algorithms).;%
Deisenroth, Faisal and Ong, \textit{Mathematics for Machine Learning}, Ch. 7.;%
Kingma and Ba, ``Adam: A Method for Stochastic Optimization,'' ICLR 2015.;%
Ruder, ``An overview of gradient descent optimization algorithms,'' arXiv:1609.04747, 2016.%
}

\newcommand{\AMLUnitIVSources}{%
Han, Kamber and Pei, \textit{Data Mining: Concepts and Techniques} (3e), Ch. 3.;%
M\"uller and Guido, \textit{Introduction to Machine Learning with Python}, Ch. 3--4.;%
James, Witten, Hastie, Tibshirani and Taylor, \textit{An Introduction to Statistical Learning with Python}, Ch. 5--6, 12.;%
Scikit-learn documentation (preprocessing, model selection), accessed 2026-08-04.%
}

\newcommand{\AMLUnitVSources}{%
James et al., \textit{An Introduction to Statistical Learning with Python}, Ch. 3--6.;%
Bishop, \textit{Pattern Recognition and Machine Learning}, Ch. 3.;%
Deisenroth, Faisal and Ong, \textit{Mathematics for Machine Learning}, Ch. 9.;%
M\"uller and Guido, \textit{Introduction to Machine Learning with Python}, \S2.3.%
}

\newcommand{\AMLUnitVISources}{%
Han, Kamber and Pei, \textit{Data Mining: Concepts and Techniques} (3e), Ch. 8--9.;%
James et al., \textit{An Introduction to Statistical Learning with Python}, Ch. 8--9.;%
Bishop, \textit{Pattern Recognition and Machine Learning}, Ch. 4--5, 7, 14.;%
Zhang et al., \textit{Dive into Deep Learning}, Ch. 5.%
}

\newcommand{\AMLUnitVIISources}{%
Han, Kamber and Pei, \textit{Data Mining: Concepts and Techniques} (3e), Ch. 10.;%
James et al., \textit{An Introduction to Statistical Learning with Python}, Ch. 12.;%
M\"uller and Guido, \textit{Introduction to Machine Learning with Python}, \S3.5.;%
Ester, Kriegel, Sander and Xu, ``A density-based algorithm for discovering clusters (DBSCAN),'' KDD 1996.%
}

\newcommand{\AMLLabSources}{%
M\"uller and Guido, \textit{Introduction to Machine Learning with Python}, companion notebooks.;%
McKinney, \textit{Python for Data Analysis} (3e), O'Reilly, 2022.;%
Scikit-learn user guide, accessed 2026-08-04.;%
Matplotlib and Seaborn documentation, accessed 2026-08-04.%
}
%
}{%
  % Faculty copies live under `private/` so their compile CWD is different.
  \IfFileExists{../../../../public/_shared/aml-sources.tex}{%
    % Source strings for Applied Machine Learning (CSAI2017P) — used by slides + notes.
% Prescribed texts and reference books are taken from the course syllabus.
% Support texts are the author-hosted open-access editions:
%   statlearning.com            (ISL with Python)
%   probml.github.io/pml-book   (Murphy, Probabilistic Machine Learning)
%   d2l.ai                      (Dive into Deep Learning)
%   mml-book.github.io          (Mathematics for Machine Learning)

\newcommand{\AMLRepoURL}{https://github.com/tali7c/Applied-Machine-Learning}

% --- prescribed -------------------------------------------------------------
\newcommand{\AMLPrescribed}{%
M\"uller and Guido, \textit{Introduction to Machine Learning with Python}, O'Reilly, 2016.;%
Bishop, \textit{Pattern Recognition and Machine Learning}, Springer, 2006.;%
Murphy, \textit{Machine Learning: A Probabilistic Perspective}, MIT Press, 2012.;%
Han, Kamber and Pei, \textit{Data Mining: Concepts and Techniques} (3e), Morgan Kaufmann, 2012.%
}

% --- unit-wise ---------------------------------------------------------------
\newcommand{\AMLUnitISources}{%
M\"uller and Guido, \textit{Introduction to Machine Learning with Python}, Ch. 1.;%
Bishop, \textit{Pattern Recognition and Machine Learning}, Ch. 1.;%
Murphy, \textit{Probabilistic Machine Learning: An Introduction}, MIT Press, 2022, Ch. 1.;%
Mitchell, \textit{Machine Learning}, McGraw Hill, 1997, Ch. 1.;%
Scikit-learn documentation, accessed 2026-08-04.%
}

\newcommand{\AMLUnitIISources}{%
Bishop, \textit{Pattern Recognition and Machine Learning}, \S1.5, \S4.3, \S7.1.;%
Zhang, Lipton, Li and Smola, \textit{Dive into Deep Learning}, \S3.4.;%
Shalev-Shwartz and Ben-David, \textit{Understanding Machine Learning}, Ch. 15.;%
Murphy, \textit{Probabilistic Machine Learning: An Introduction}, Ch. 5.%
}

\newcommand{\AMLUnitIIISources}{%
Zhang, Lipton, Li and Smola, \textit{Dive into Deep Learning}, Ch. 12 (Optimization Algorithms).;%
Deisenroth, Faisal and Ong, \textit{Mathematics for Machine Learning}, Ch. 7.;%
Kingma and Ba, ``Adam: A Method for Stochastic Optimization,'' ICLR 2015.;%
Ruder, ``An overview of gradient descent optimization algorithms,'' arXiv:1609.04747, 2016.%
}

\newcommand{\AMLUnitIVSources}{%
Han, Kamber and Pei, \textit{Data Mining: Concepts and Techniques} (3e), Ch. 3.;%
M\"uller and Guido, \textit{Introduction to Machine Learning with Python}, Ch. 3--4.;%
James, Witten, Hastie, Tibshirani and Taylor, \textit{An Introduction to Statistical Learning with Python}, Ch. 5--6, 12.;%
Scikit-learn documentation (preprocessing, model selection), accessed 2026-08-04.%
}

\newcommand{\AMLUnitVSources}{%
James et al., \textit{An Introduction to Statistical Learning with Python}, Ch. 3--6.;%
Bishop, \textit{Pattern Recognition and Machine Learning}, Ch. 3.;%
Deisenroth, Faisal and Ong, \textit{Mathematics for Machine Learning}, Ch. 9.;%
M\"uller and Guido, \textit{Introduction to Machine Learning with Python}, \S2.3.%
}

\newcommand{\AMLUnitVISources}{%
Han, Kamber and Pei, \textit{Data Mining: Concepts and Techniques} (3e), Ch. 8--9.;%
James et al., \textit{An Introduction to Statistical Learning with Python}, Ch. 8--9.;%
Bishop, \textit{Pattern Recognition and Machine Learning}, Ch. 4--5, 7, 14.;%
Zhang et al., \textit{Dive into Deep Learning}, Ch. 5.%
}

\newcommand{\AMLUnitVIISources}{%
Han, Kamber and Pei, \textit{Data Mining: Concepts and Techniques} (3e), Ch. 10.;%
James et al., \textit{An Introduction to Statistical Learning with Python}, Ch. 12.;%
M\"uller and Guido, \textit{Introduction to Machine Learning with Python}, \S3.5.;%
Ester, Kriegel, Sander and Xu, ``A density-based algorithm for discovering clusters (DBSCAN),'' KDD 1996.%
}

\newcommand{\AMLLabSources}{%
M\"uller and Guido, \textit{Introduction to Machine Learning with Python}, companion notebooks.;%
McKinney, \textit{Python for Data Analysis} (3e), O'Reilly, 2022.;%
Scikit-learn user guide, accessed 2026-08-04.;%
Matplotlib and Seaborn documentation, accessed 2026-08-04.%
}
%
  }{%
    % Question papers live at `private/<family>/<instrument>/`.
    \IfFileExists{../../../public/_shared/aml-sources.tex}{%
      % Source strings for Applied Machine Learning (CSAI2017P) — used by slides + notes.
% Prescribed texts and reference books are taken from the course syllabus.
% Support texts are the author-hosted open-access editions:
%   statlearning.com            (ISL with Python)
%   probml.github.io/pml-book   (Murphy, Probabilistic Machine Learning)
%   d2l.ai                      (Dive into Deep Learning)
%   mml-book.github.io          (Mathematics for Machine Learning)

\newcommand{\AMLRepoURL}{https://github.com/tali7c/Applied-Machine-Learning}

% --- prescribed -------------------------------------------------------------
\newcommand{\AMLPrescribed}{%
M\"uller and Guido, \textit{Introduction to Machine Learning with Python}, O'Reilly, 2016.;%
Bishop, \textit{Pattern Recognition and Machine Learning}, Springer, 2006.;%
Murphy, \textit{Machine Learning: A Probabilistic Perspective}, MIT Press, 2012.;%
Han, Kamber and Pei, \textit{Data Mining: Concepts and Techniques} (3e), Morgan Kaufmann, 2012.%
}

% --- unit-wise ---------------------------------------------------------------
\newcommand{\AMLUnitISources}{%
M\"uller and Guido, \textit{Introduction to Machine Learning with Python}, Ch. 1.;%
Bishop, \textit{Pattern Recognition and Machine Learning}, Ch. 1.;%
Murphy, \textit{Probabilistic Machine Learning: An Introduction}, MIT Press, 2022, Ch. 1.;%
Mitchell, \textit{Machine Learning}, McGraw Hill, 1997, Ch. 1.;%
Scikit-learn documentation, accessed 2026-08-04.%
}

\newcommand{\AMLUnitIISources}{%
Bishop, \textit{Pattern Recognition and Machine Learning}, \S1.5, \S4.3, \S7.1.;%
Zhang, Lipton, Li and Smola, \textit{Dive into Deep Learning}, \S3.4.;%
Shalev-Shwartz and Ben-David, \textit{Understanding Machine Learning}, Ch. 15.;%
Murphy, \textit{Probabilistic Machine Learning: An Introduction}, Ch. 5.%
}

\newcommand{\AMLUnitIIISources}{%
Zhang, Lipton, Li and Smola, \textit{Dive into Deep Learning}, Ch. 12 (Optimization Algorithms).;%
Deisenroth, Faisal and Ong, \textit{Mathematics for Machine Learning}, Ch. 7.;%
Kingma and Ba, ``Adam: A Method for Stochastic Optimization,'' ICLR 2015.;%
Ruder, ``An overview of gradient descent optimization algorithms,'' arXiv:1609.04747, 2016.%
}

\newcommand{\AMLUnitIVSources}{%
Han, Kamber and Pei, \textit{Data Mining: Concepts and Techniques} (3e), Ch. 3.;%
M\"uller and Guido, \textit{Introduction to Machine Learning with Python}, Ch. 3--4.;%
James, Witten, Hastie, Tibshirani and Taylor, \textit{An Introduction to Statistical Learning with Python}, Ch. 5--6, 12.;%
Scikit-learn documentation (preprocessing, model selection), accessed 2026-08-04.%
}

\newcommand{\AMLUnitVSources}{%
James et al., \textit{An Introduction to Statistical Learning with Python}, Ch. 3--6.;%
Bishop, \textit{Pattern Recognition and Machine Learning}, Ch. 3.;%
Deisenroth, Faisal and Ong, \textit{Mathematics for Machine Learning}, Ch. 9.;%
M\"uller and Guido, \textit{Introduction to Machine Learning with Python}, \S2.3.%
}

\newcommand{\AMLUnitVISources}{%
Han, Kamber and Pei, \textit{Data Mining: Concepts and Techniques} (3e), Ch. 8--9.;%
James et al., \textit{An Introduction to Statistical Learning with Python}, Ch. 8--9.;%
Bishop, \textit{Pattern Recognition and Machine Learning}, Ch. 4--5, 7, 14.;%
Zhang et al., \textit{Dive into Deep Learning}, Ch. 5.%
}

\newcommand{\AMLUnitVIISources}{%
Han, Kamber and Pei, \textit{Data Mining: Concepts and Techniques} (3e), Ch. 10.;%
James et al., \textit{An Introduction to Statistical Learning with Python}, Ch. 12.;%
M\"uller and Guido, \textit{Introduction to Machine Learning with Python}, \S3.5.;%
Ester, Kriegel, Sander and Xu, ``A density-based algorithm for discovering clusters (DBSCAN),'' KDD 1996.%
}

\newcommand{\AMLLabSources}{%
M\"uller and Guido, \textit{Introduction to Machine Learning with Python}, companion notebooks.;%
McKinney, \textit{Python for Data Analysis} (3e), O'Reilly, 2022.;%
Scikit-learn user guide, accessed 2026-08-04.;%
Matplotlib and Seaborn documentation, accessed 2026-08-04.%
}
%
    }{%
      % Marking schemes live at `private/solutions/<family>/<id>/latex/`.
      \IfFileExists{../../../../../public/_shared/aml-sources.tex}{%
        % Source strings for Applied Machine Learning (CSAI2017P) — used by slides + notes.
% Prescribed texts and reference books are taken from the course syllabus.
% Support texts are the author-hosted open-access editions:
%   statlearning.com            (ISL with Python)
%   probml.github.io/pml-book   (Murphy, Probabilistic Machine Learning)
%   d2l.ai                      (Dive into Deep Learning)
%   mml-book.github.io          (Mathematics for Machine Learning)

\newcommand{\AMLRepoURL}{https://github.com/tali7c/Applied-Machine-Learning}

% --- prescribed -------------------------------------------------------------
\newcommand{\AMLPrescribed}{%
M\"uller and Guido, \textit{Introduction to Machine Learning with Python}, O'Reilly, 2016.;%
Bishop, \textit{Pattern Recognition and Machine Learning}, Springer, 2006.;%
Murphy, \textit{Machine Learning: A Probabilistic Perspective}, MIT Press, 2012.;%
Han, Kamber and Pei, \textit{Data Mining: Concepts and Techniques} (3e), Morgan Kaufmann, 2012.%
}

% --- unit-wise ---------------------------------------------------------------
\newcommand{\AMLUnitISources}{%
M\"uller and Guido, \textit{Introduction to Machine Learning with Python}, Ch. 1.;%
Bishop, \textit{Pattern Recognition and Machine Learning}, Ch. 1.;%
Murphy, \textit{Probabilistic Machine Learning: An Introduction}, MIT Press, 2022, Ch. 1.;%
Mitchell, \textit{Machine Learning}, McGraw Hill, 1997, Ch. 1.;%
Scikit-learn documentation, accessed 2026-08-04.%
}

\newcommand{\AMLUnitIISources}{%
Bishop, \textit{Pattern Recognition and Machine Learning}, \S1.5, \S4.3, \S7.1.;%
Zhang, Lipton, Li and Smola, \textit{Dive into Deep Learning}, \S3.4.;%
Shalev-Shwartz and Ben-David, \textit{Understanding Machine Learning}, Ch. 15.;%
Murphy, \textit{Probabilistic Machine Learning: An Introduction}, Ch. 5.%
}

\newcommand{\AMLUnitIIISources}{%
Zhang, Lipton, Li and Smola, \textit{Dive into Deep Learning}, Ch. 12 (Optimization Algorithms).;%
Deisenroth, Faisal and Ong, \textit{Mathematics for Machine Learning}, Ch. 7.;%
Kingma and Ba, ``Adam: A Method for Stochastic Optimization,'' ICLR 2015.;%
Ruder, ``An overview of gradient descent optimization algorithms,'' arXiv:1609.04747, 2016.%
}

\newcommand{\AMLUnitIVSources}{%
Han, Kamber and Pei, \textit{Data Mining: Concepts and Techniques} (3e), Ch. 3.;%
M\"uller and Guido, \textit{Introduction to Machine Learning with Python}, Ch. 3--4.;%
James, Witten, Hastie, Tibshirani and Taylor, \textit{An Introduction to Statistical Learning with Python}, Ch. 5--6, 12.;%
Scikit-learn documentation (preprocessing, model selection), accessed 2026-08-04.%
}

\newcommand{\AMLUnitVSources}{%
James et al., \textit{An Introduction to Statistical Learning with Python}, Ch. 3--6.;%
Bishop, \textit{Pattern Recognition and Machine Learning}, Ch. 3.;%
Deisenroth, Faisal and Ong, \textit{Mathematics for Machine Learning}, Ch. 9.;%
M\"uller and Guido, \textit{Introduction to Machine Learning with Python}, \S2.3.%
}

\newcommand{\AMLUnitVISources}{%
Han, Kamber and Pei, \textit{Data Mining: Concepts and Techniques} (3e), Ch. 8--9.;%
James et al., \textit{An Introduction to Statistical Learning with Python}, Ch. 8--9.;%
Bishop, \textit{Pattern Recognition and Machine Learning}, Ch. 4--5, 7, 14.;%
Zhang et al., \textit{Dive into Deep Learning}, Ch. 5.%
}

\newcommand{\AMLUnitVIISources}{%
Han, Kamber and Pei, \textit{Data Mining: Concepts and Techniques} (3e), Ch. 10.;%
James et al., \textit{An Introduction to Statistical Learning with Python}, Ch. 12.;%
M\"uller and Guido, \textit{Introduction to Machine Learning with Python}, \S3.5.;%
Ester, Kriegel, Sander and Xu, ``A density-based algorithm for discovering clusters (DBSCAN),'' KDD 1996.%
}

\newcommand{\AMLLabSources}{%
M\"uller and Guido, \textit{Introduction to Machine Learning with Python}, companion notebooks.;%
McKinney, \textit{Python for Data Analysis} (3e), O'Reilly, 2022.;%
Scikit-learn user guide, accessed 2026-08-04.;%
Matplotlib and Seaborn documentation, accessed 2026-08-04.%
}
%
      }{%
        % Source strings for Applied Machine Learning (CSAI2017P) — used by slides + notes.
% Prescribed texts and reference books are taken from the course syllabus.
% Support texts are the author-hosted open-access editions:
%   statlearning.com            (ISL with Python)
%   probml.github.io/pml-book   (Murphy, Probabilistic Machine Learning)
%   d2l.ai                      (Dive into Deep Learning)
%   mml-book.github.io          (Mathematics for Machine Learning)

\newcommand{\AMLRepoURL}{https://github.com/tali7c/Applied-Machine-Learning}

% --- prescribed -------------------------------------------------------------
\newcommand{\AMLPrescribed}{%
M\"uller and Guido, \textit{Introduction to Machine Learning with Python}, O'Reilly, 2016.;%
Bishop, \textit{Pattern Recognition and Machine Learning}, Springer, 2006.;%
Murphy, \textit{Machine Learning: A Probabilistic Perspective}, MIT Press, 2012.;%
Han, Kamber and Pei, \textit{Data Mining: Concepts and Techniques} (3e), Morgan Kaufmann, 2012.%
}

% --- unit-wise ---------------------------------------------------------------
\newcommand{\AMLUnitISources}{%
M\"uller and Guido, \textit{Introduction to Machine Learning with Python}, Ch. 1.;%
Bishop, \textit{Pattern Recognition and Machine Learning}, Ch. 1.;%
Murphy, \textit{Probabilistic Machine Learning: An Introduction}, MIT Press, 2022, Ch. 1.;%
Mitchell, \textit{Machine Learning}, McGraw Hill, 1997, Ch. 1.;%
Scikit-learn documentation, accessed 2026-08-04.%
}

\newcommand{\AMLUnitIISources}{%
Bishop, \textit{Pattern Recognition and Machine Learning}, \S1.5, \S4.3, \S7.1.;%
Zhang, Lipton, Li and Smola, \textit{Dive into Deep Learning}, \S3.4.;%
Shalev-Shwartz and Ben-David, \textit{Understanding Machine Learning}, Ch. 15.;%
Murphy, \textit{Probabilistic Machine Learning: An Introduction}, Ch. 5.%
}

\newcommand{\AMLUnitIIISources}{%
Zhang, Lipton, Li and Smola, \textit{Dive into Deep Learning}, Ch. 12 (Optimization Algorithms).;%
Deisenroth, Faisal and Ong, \textit{Mathematics for Machine Learning}, Ch. 7.;%
Kingma and Ba, ``Adam: A Method for Stochastic Optimization,'' ICLR 2015.;%
Ruder, ``An overview of gradient descent optimization algorithms,'' arXiv:1609.04747, 2016.%
}

\newcommand{\AMLUnitIVSources}{%
Han, Kamber and Pei, \textit{Data Mining: Concepts and Techniques} (3e), Ch. 3.;%
M\"uller and Guido, \textit{Introduction to Machine Learning with Python}, Ch. 3--4.;%
James, Witten, Hastie, Tibshirani and Taylor, \textit{An Introduction to Statistical Learning with Python}, Ch. 5--6, 12.;%
Scikit-learn documentation (preprocessing, model selection), accessed 2026-08-04.%
}

\newcommand{\AMLUnitVSources}{%
James et al., \textit{An Introduction to Statistical Learning with Python}, Ch. 3--6.;%
Bishop, \textit{Pattern Recognition and Machine Learning}, Ch. 3.;%
Deisenroth, Faisal and Ong, \textit{Mathematics for Machine Learning}, Ch. 9.;%
M\"uller and Guido, \textit{Introduction to Machine Learning with Python}, \S2.3.%
}

\newcommand{\AMLUnitVISources}{%
Han, Kamber and Pei, \textit{Data Mining: Concepts and Techniques} (3e), Ch. 8--9.;%
James et al., \textit{An Introduction to Statistical Learning with Python}, Ch. 8--9.;%
Bishop, \textit{Pattern Recognition and Machine Learning}, Ch. 4--5, 7, 14.;%
Zhang et al., \textit{Dive into Deep Learning}, Ch. 5.%
}

\newcommand{\AMLUnitVIISources}{%
Han, Kamber and Pei, \textit{Data Mining: Concepts and Techniques} (3e), Ch. 10.;%
James et al., \textit{An Introduction to Statistical Learning with Python}, Ch. 12.;%
M\"uller and Guido, \textit{Introduction to Machine Learning with Python}, \S3.5.;%
Ester, Kriegel, Sander and Xu, ``A density-based algorithm for discovering clusters (DBSCAN),'' KDD 1996.%
}

\newcommand{\AMLLabSources}{%
M\"uller and Guido, \textit{Introduction to Machine Learning with Python}, companion notebooks.;%
McKinney, \textit{Python for Data Analysis} (3e), O'Reilly, 2022.;%
Scikit-learn user guide, accessed 2026-08-04.;%
Matplotlib and Seaborn documentation, accessed 2026-08-04.%
}
%
      }%
    }%
  }%
}


\usepackage{longtable}

\title{Applied Machine Learning (CSAI2017P)\\[4pt]
       \large Student Course Handbook}
\author{Dr.\ Tofik Ali \\ School of Computer Science, UPES Dehradun}
\date{B.Tech CSE (AI \& ML), Semester IV \quad\textbullet\quad Autumn 2026}

\begin{document}
\maketitle

\begin{keyidea}[What this handbook is]
Everything you need to know about how this course runs: what you will learn,
how you will be assessed, when each thing is due, and where to get help.

Keep it. Almost every question you are likely to ask this semester is answered
somewhere in these pages. Anything that changes will be announced on the LMS and
this document reissued.
\end{keyidea}

\tableofcontents
\newpage

% =========================================================================
\section{The course at a glance}

\begin{center}
\begin{tabular}{@{}ll@{}}
\toprule
Course title & Applied Machine Learning \\
Course code & CSAI2017 (theory) \textbullet\ \textbf{CSAI2017P} (theory + lab) \\
Credits (L-T-P-C) & 4-0-1-5 \\
Programme & B.Tech CSE (Artificial Intelligence \& Machine Learning) \\
Semester & IV \\
Session & Autumn 2026, AY 2026-27 (confirm on the timetable) \\
Units & 7 \\
Theory & 24 lecture packages across 40 sessions \\
Laboratory & 15 experiments \\
Prerequisites & Elements of AI \& ML (CSAI2018); Python Programming (CSEG1021) \\
Faculty & Dr.\ Tofik Ali \\
Repository & \url{https://github.com/tali7c/Applied-Machine-Learning} \\
\bottomrule
\end{tabular}
\end{center}

\subsection{Course outcomes}

By the end of the course you should be able to:

\begin{itemize}[leftmargin=*]
  \item \textbf{CO1} --- Understand the core concepts and techniques of machine
    learning and artificial intelligence.
  \item \textbf{CO2} --- Develop machine learning models using popular libraries
    and frameworks.
  \item \textbf{CO3} --- Evaluate the performance of machine learning models
    using appropriate metrics.
  \item \textbf{CO4} --- Apply machine learning to various real-world problems
    and domains.
\end{itemize}

\subsection{What you will actually be able to do}

By the end of the semester you will have built, from scratch, a housing-price
predictor, a stock forecaster, a churn model, a digit recogniser, a spam filter,
a credit scorer, an anomaly detector, a breast-cancer diagnostic model and a
crop-disease classifier --- and, more importantly, you will be able to tell when
one of them is lying to you.

\newpage
% =========================================================================
\section{Syllabus}

The order is deliberate. Units II and III come early because once you know
\textbf{what quantity we are minimising} and \textbf{how the minimisation is
carried out}, every method in Units V--VII is a variation on one theme rather
than a new thing to memorise.

\begin{longtable}{@{}p{0.6cm}p{3.2cm}p{1.7cm}p{8.9cm}@{}}
\toprule
\textbf{} & \textbf{Unit} & \textbf{Sessions} & \textbf{Content}\\
\midrule
\endhead
I & Introduction & 2 &
Overview of machine learning and its applications. Types of machine learning:
supervised, unsupervised, reinforcement. Python and its libraries for ML
(NumPy, pandas, scikit-learn).\\[4pt]

II & Loss Functions & 2 &
Mean Squared Error, Mean Absolute Error, Huber loss. Binary, Categorical and
Sparse Categorical Cross-Entropy. Hinge loss (SVM loss), Triplet loss.\\[4pt]

III & Optimizer Functions & 4 &
Stochastic gradient descent. Mini-batch gradient descent. Momentum. Adagrad.
Adam. RMSprop. Adadelta. Numerical practice.\\[4pt]

IV & Data Preprocessing & 6 &
Data cleaning; handling missing data and outliers; data transformation. Feature
engineering, scaling and encoding. Data reduction and dimensionality reduction;
feature selection. Train--test splitting and cross-validation. Handling
imbalanced data.\\[4pt]

V & Regression & 7 &
Introduction to regression. Simple linear regression and its proof. Least-squares
estimation; line of best fit. Maximum likelihood estimation. $R^2$; model
adequacy; over-fitting and its detection. Logistic regression. Multiple linear
regression and coefficient interpretation. Ridge and lasso regularization.\\[4pt]

VI & Classification & 12 &
ML classifiers and instance-based learning. K-Nearest Neighbour. Decision trees;
information gain; ID3. Bayesian algorithms. Ensembles: bagging, boosting, random
forests. Neural networks; activation functions; multi-layer perceptron;
backpropagation. Support vector machines. Model evaluation; ROC and AUC.\\[4pt]

VII & Clustering Techniques & 7 &
Introduction to clustering; cluster statistics and applications. Similarity and
dissimilarity; data types in clustering. Hierarchical clustering and HAC;
linkage methods; cluster distance measures. K-means; K-medoids; CLARA.
Density-based methods; DBSCAN.\\
\bottomrule
\end{longtable}

\noindent Total: \textbf{40 theory sessions}, including two assessment sessions
and a final revision session.

\begin{caution}[About the session counts]
The unit contents above are the official syllabus. The \emph{per-unit session
counts} are this course's planned allocation of the available contact sessions,
not a figure from the Academic Handbook, and they may shift by a session or two
as the semester runs. Treat them as a guide to the relative weight of each unit,
not as a contract.
\end{caution}

\newpage
% =========================================================================
\section{How you are assessed}

\begin{caution}[Indicative, pending approval]
The component split below is subject to departmental approval. The confirmed
scheme will be posted on the LMS, and this handbook reissued if it changes. The
Internal 50 : Mid-Semester 20 : End-Semester 30 frame is likewise to be confirmed
against the Academic Handbook.
\end{caution}

\subsection{The 100 marks}

\begin{center}
\begin{tabular}{@{}p{6.2cm}cp{6.4cm}@{}}
\toprule
\textbf{Component} & \textbf{Marks} & \textbf{What it is}\\
\midrule
Theory assignments & 5 & Two assignments, each followed by a short in-class assessment test\\
Theory quizzes / tests & 10 & Ten 20-minute tests through the semester\\
Project --- individual & 10 & A quiz on your own group's project; \textbf{marked individually}\\
Project --- group & 5 & The state of what your group submitted\\
Lab assignments & 10 & Ten assignments covering experiments 1--14\\
Lab quizzes & 10 & Ten short quizzes on \emph{your own} submitted lab code\\
Mid-Semester examination & 20 & Common paper, university schedule\\
End-Semester examination & 30 & Common paper, university schedule\\
\midrule
\textbf{Total} & \textbf{100} & \\
\bottomrule
\end{tabular}
\end{center}

\subsection{How that maps to the university scheme}

\begin{center}
\begin{tabular}{@{}lcl@{}}
\toprule
\textbf{Block} & \textbf{Marks} & \textbf{Made up of}\\
\midrule
Continuous Internal Assessment (CIA) & 50 & everything except the two examinations\\
Mid-Semester Examination & 20 & \\
End-Semester Examination & 30 & \\
\bottomrule
\end{tabular}
\end{center}

\begin{keyidea}[The single most important line in this handbook]
\textbf{Seventy of your hundred marks are decided before the end-semester
examination} --- fifty from work spread across the whole semester (quizzes,
assignments, lab work, project) and twenty at the mid-semester. A student who
coasts until the last month cannot recover, because the marks are already gone.
\end{keyidea}

\subsection{Grades}

\begin{center}
\begin{tabular}{@{}lcccccccc@{}}
\toprule
\textbf{Total} & $\geq 90$ & $\geq 80$ & $\geq 70$ & $\geq 60$ & $\geq 50$ & $\geq 45$ & $\geq 40$ & $< 40$\\
\midrule
\textbf{Grade} & O & A+ & A & B+ & B & C & P & F\\
\bottomrule
\end{tabular}
\end{center}

\newpage
% =========================================================================
\section{The assessments in detail}

\subsection{Theory quizzes --- 10 marks}

Ten tests, T1--T10. Each is a \textbf{20-minute} instrument run at the close of
a lecture session; they do not consume a whole session. The diagnostic Test-0 is
the exception at \textbf{30 minutes}.

\begin{center}
\begin{tabular}{@{}llp{8.4cm}@{}}
\toprule
\textbf{Test} & \textbf{Session} & \textbf{Coverage}\\
\midrule
T0 & week 2, TBC & Prerequisites and general understanding --- \textbf{no marks}\\
T1 & 6 & Units I--II\\
T2 & 10 & Unit III; Unit IV (cleaning)\\
T3 & 13 & Unit IV\\
T4 & 17 & Unit V --- foundations, least squares, MLE\\
T5 & 21 & Unit V --- logistic, MLR, regularization\\
T6 & 25 & Unit VI --- foundations, KNN, decision trees\\
T7 & 28 & Unit VI --- Bayesian, ensembles\\
T8 & 32 & Unit VI --- neural networks, SVM, evaluation\\
T9 & 35 & Unit VII --- foundations, similarity, data types\\
T10 & 39 & Unit VII --- hierarchical, partitioning, density-based\\
\bottomrule
\end{tabular}
\end{center}

Each test is marked out of 10; the component is the total scaled to 10.
\textbf{An absence counts as zero}, so tell me in advance if you cannot sit one.

A \textbf{drop-lowest-two} provision may be applied, in which case your best 8 of
10 count. Whether it is in force will be confirmed early in the semester ---
plan on all ten counting and treat it as a safety net, not a budget.

\subsubsection*{About Test-0, and about the different question sets}

\textbf{Test-0} is a diagnostic. It carries \textbf{no marks}, is not a ranking,
and is not an ability test. There is nothing to revise for it. Its purpose is to
let the teaching team see \emph{how you approach a problem}, so that the short
tests through the semester can be framed in a way that suits how you work. Show
your working; an incomplete answer that makes its reasoning visible is more
useful than a bare correct one.

You will notice that the quizzes come in several versions. What you should know
about that:

\begin{itemize}[leftmargin=*]
  \item The versions match the \textbf{style} of a question to how you work.
    They are not ability tiers.
  \item Every version asks \textbf{equally demanding} questions, on the same
    topics, for the same marks, and is marked against the \textbf{same rubric}.
  \item Which version you get may change during the semester. That is normal and
    is not a promotion or a demotion.
  \item It has \textbf{no effect whatsoever} on any mark you receive.
  \item T1 may run as a single common paper for everyone; the versions begin at
    T2 in that case.
\end{itemize}

If any of that does not match your experience, tell me --- that would be a fault
in the design and I would want to fix it.

\subsection{Theory assignments --- 5 marks}

Two assignments, each worth 2.5. Each has two parts: the submission itself
(1.0) and a short \textbf{in-class assessment test} on what you submitted (1.5).

\begin{center}
\begin{tabular}{@{}lp{6.6cm}ll@{}}
\toprule
\textbf{Asst} & \textbf{Coverage} & \textbf{Issued} & \textbf{Submission + test}\\
\midrule
A1 & Units I--IV --- preprocessing and cross-validation pipeline & after session 13 & \textbf{session 14}\\
A2 & Units V--VI --- classifier comparison with ROC/AUC & after session 32 & \textbf{session 33}\\
\bottomrule
\end{tabular}
\end{center}

\begin{caution}[Why the test is worth more than the submission]
Because it is the part that proves the work is yours. If you did the assignment,
the test is straightforward. If you did not, no amount of polish on the
submission will save it.
\end{caution}

\subsection{Laboratory --- 20 marks}

Experiments 1--14 are grouped into \textbf{ten assignments}. Experiment 15, the
mini-project, is assessed under the project component instead.

\begin{center}
\begin{tabular}{@{}llp{5.5cm}cc@{}}
\toprule
\textbf{Asst} & \textbf{Expts} & \textbf{Topic} & \textbf{Due after} & \textbf{Quiz}\\
 & & & \textbf{lab} & \\
\midrule
LA1 & 1, 2 & Environment setup; preprocessing and EDA & 2 & LQ1, lab 3\\
LA2 & 3 & Regression: housing price prediction & 3 & LQ2, lab 4\\
LA3 & 4 & Time-series regression: stock prices & 4 & LQ3, lab 5\\
LA4 & 5 & Customer churn prediction & 5 & LQ4, lab 6\\
LA5 & 6, 7 & Image and multiclass classification: digits and Iris & 7 & LQ5, lab 8\\
LA6 & 8 & Spam email detection & 8 & LQ6, lab 9\\
LA7 & 9 & Credit risk assessment & 9 & LQ7, lab 10\\
LA8 & 10, 11 & Anomaly detection; physiological signal classification & 11 & LQ8, lab 12\\
LA9 & 12 & Breast cancer diagnosis from imaging features & 12 & LQ9, lab 13\\
LA10 & 13, 14 & Clustering medical images; crop disease classification & 14 & LQ10, lab 15\\
\bottomrule
\end{tabular}
\end{center}

\subsubsection*{How a lab assignment is marked}

Each is out of 10, with a choice structure:

\begin{center}
\begin{tabular}{@{}llll@{}}
\toprule
\textbf{Section} & \textbf{Level} & \textbf{Choice rule} & \textbf{Marks}\\
\midrule
A & Easy & attempt any 2 of 3 (2 marks each) & 4\\
B & Medium & attempt any 1 of 2 (4 marks each) & 4\\
C & Consolidation & compulsory & 2\\
\midrule
\multicolumn{3}{@{}l}{\textbf{Total}} & \textbf{10}\\
\bottomrule
\end{tabular}
\end{center}

Attempting more than the rule allows is not penalised --- the best attempts are
marked.

\subsubsection*{Lab quizzes --- and why you should write your own code}

Each lab quiz asks you to \textbf{read, trace, debug and modify your own
submitted code}. Not recall, not theory --- your code, in front of you, with
questions about it.

\begin{caution}
This is deliberate, and it is the reason copying does not work in this course.
Code you did not write is very hard to explain on the spot. Struggling through
your own imperfect solution will score better than submitting someone else's
good one.
\end{caution}

\subsection{Project --- 15 marks}

Ten projects are offered. You form your own group and choose one. Each project is
normally taken by one group only; two groups may take the same one only with
prior approval.

\begin{center}
\begin{tabular}{@{}p{5.4cm}cp{6.8cm}@{}}
\toprule
\textbf{Sub-component} & \textbf{Marks} & \textbf{Basis}\\
\midrule
Project quiz & 10 & An oral and written test on your group's project. \textbf{Marked individually} --- everyone is asked separately.\\
Project state & 5 & Completeness, correctness, reproducibility and documentation of the artefact. \textbf{Awarded to the group.}\\
\bottomrule
\end{tabular}
\end{center}

\begin{center}
\begin{tabular}{@{}lc@{}}
\toprule
\textbf{Milestone} & \textbf{Session}\\
\midrule
Project catalogue released & 8\\
Groups formed and project selected & 13\\
Checkpoint 1 --- data secured, baseline model & 21\\
Checkpoint 2 --- full pipeline, initial results & 32\\
Final submission (code, report, reproducibility notes) & 39\\
Individual project quiz & 40 / lab 15\\
\bottomrule
\end{tabular}
\end{center}

\begin{keyidea}[Read this before you choose your group]
Two thirds of the project mark is \textbf{individual}. A strong contributor in a
weak group can still score well. A passenger in a strong group cannot --- the
group mark will not save them, because it is only 5 of the 15.

Group marks are not divided by contribution. All individual differentiation
happens through the individual quiz.
\end{keyidea}

\newpage
% =========================================================================
\section{Semester planner}

Sessions are numbered rather than dated; dates follow the timetable and will be
announced. Print this page.

\begin{center}
\begin{tabular}{@{}clp{7.2cm}@{}}
\toprule
\textbf{Session} & \textbf{What happens} & \textbf{What you should have done}\\
\midrule
1--2 & Unit I lectures & \texttt{pip install numpy pandas scikit-learn matplotlib seaborn jupyter}\\
week 2 & \textbf{Test-0} (no marks) & Nothing --- turn up\\
lab 2 & \textbf{LA1 due} & Experiments 1--2 written up\\
lab 3 & \textbf{LQ1} & Be able to explain your LA1 code\\
6 & \textbf{T1} & Units I--II\\
8 & Project catalogue released & Start thinking about a group\\
10 & \textbf{T2} & Unit III; Unit IV cleaning\\
13 & \textbf{T3}; groups and project fixed & Unit IV\\
14 & \textbf{A1 submission + assessment test} & Assignment 1 done and understood\\
17 & \textbf{T4} & Unit V foundations\\
21 & \textbf{T5}; project checkpoint 1 & Unit V; data secured, baseline model\\
25 & \textbf{T6} & Unit VI foundations, KNN, trees\\
28 & \textbf{T7} & Unit VI Bayesian, ensembles\\
32 & \textbf{T8}; project checkpoint 2 & Unit VI complete; pipeline running\\
33 & \textbf{A2 submission + assessment test} & Assignment 2 done and understood\\
35 & \textbf{T9} & Unit VII foundations\\
39 & \textbf{T10}; project final submission & Unit VII complete; project handed in\\
40 & Revision; project quiz & Everything\\
\bottomrule
\end{tabular}
\end{center}

\noindent Mid-Semester and End-Semester examinations follow the university
calendar and are notified separately.

\newpage
% =========================================================================
\section{How to do well in this course}

\subsection{Three rules that apply to everything}

\begin{enumerate}[leftmargin=*]
  \item \textbf{Fit every transformer inside a \texttt{Pipeline}.} Never fit
    anything on the test set --- not a scaler, not an imputer, not a feature
    selector. This is the single most common way to produce a wonderful result
    that means nothing.
  \item \textbf{Set \texttt{random\_state} wherever reproducibility is asked
    for.} A result you cannot reproduce is a result you cannot defend.
  \item \textbf{Every question wants a short written observation}, not just a
    number. Code with no commentary is capped at half marks.
\end{enumerate}

\subsection{Practical advice}

\begin{itemize}[leftmargin=*]
  \item \textbf{Restart the kernel and run top to bottom before you submit.} Out
    of order cells are the commonest reason a submission does not reproduce.
  \item \textbf{Download each dataset once} into a local \texttt{data/} folder
    and keep it for the semester. \texttt{Anchor-Datasets.pdf} has the links and
    the loading code. Some loaders need internet on first use --- run them at
    home, not in the lab.
  \item \textbf{Report metrics with units.} ``MAE 0.53'' means nothing;
    ``MAE 0.53, i.e.\ about \$53{,}000'' means something.
  \item \textbf{Do not quote more precision than you have.} If cross-validation
    gives $0.953 \pm 0.052$, writing ``95.33\%'' is false precision.
  \item \textbf{Ask in the lab rather than guessing.} That is what the session
    is for, and it costs you nothing.
\end{itemize}

\subsection{What to submit}

\begin{itemize}[leftmargin=*]
  \item Notebook: \texttt{AML\_LAxx\_<SAPID>.ipynb}, all cells executed with
    outputs visible.
  \item Report: 2--3 pages PDF, or a \texttt{README.md}, carrying your results
    table and observations.
  \item Do not upload datasets. Cite the source and include the loading code.
\end{itemize}

\subsection{Academic integrity}

Discussing ideas with classmates is encouraged. Submitting someone else's code
or text as your own is not, and the assessment design will find it: the
assignment assessment tests and the lab quizzes both ask you to explain work you
have submitted. If you used a tool or a tutorial substantially, say so in your
report --- that is honest practice, not an admission of weakness.

\newpage
% =========================================================================
\section{Books, courses and other help}

All links checked 4 August 2026. Free and no-login resources are marked
\textbf{[free]}.

\subsection*{Prescribed}

\begin{itemize}[leftmargin=*]
  \item M\"uller, A.\ C.\ and Guido, S. \emph{Introduction to Machine Learning
    with Python}. O'Reilly, 2016.
  \item Bishop, C.\ M. \emph{Pattern Recognition and Machine Learning}.
    Springer, 2006. \textbf{[free]} from Microsoft Research.
\end{itemize}

\subsection*{Reference}

\begin{itemize}[leftmargin=*]
  \item Murphy, K.\ P. \emph{Machine Learning: A Probabilistic Perspective}.
    MIT Press, 2012. The author's free successor volume is at
    \textbf{[free]} \url{https://probml.github.io/pml-book/book1.html}
  \item Han, J., Kamber, M.\ and Pei, J. \emph{Data Mining: Concepts and
    Techniques}, 3rd ed. Morgan Kaufmann, 2012. Especially Chapters 3, 8 and 10.
  \item \emph{scikit-learn user guide} --- the syllabus web resource.
    \textbf{[free]} \url{https://scikit-learn.org/stable/}
\end{itemize}

\subsection*{Open-access support (not on the official list, but good)}

\begin{itemize}[leftmargin=*]
  \item James, G.\ et al. \emph{An Introduction to Statistical Learning with
    Python}. \textbf{[free]} \url{https://www.statlearning.com/} --- the friendliest
    treatment of regression, regularization and cross-validation available.
\end{itemize}

\subsection*{Courses --- free, and worth your time}

\begin{itemize}[leftmargin=*]
  \item \textbf{NPTEL, Introduction to Machine Learning}, Prof.\ Balaraman
    Ravindran, IIT Madras. The closest match to this syllabus.
    \textbf{[free]} \url{https://nptel.ac.in/courses/106106139}
  \item \textbf{SWAYAM} --- national portal hosting NPTEL and other Indian
    university courses. \textbf{[free]} \url{https://swayam.gov.in/}
  \item \textbf{Google Machine Learning Crash Course}. \textbf{[free]}
    \url{https://developers.google.com/machine-learning/crash-course}
  \item \textbf{Kaggle Learn, Intro to Machine Learning} --- four hours, runs in
    the browser, no installation. \textbf{[free]}
    \url{https://www.kaggle.com/learn/intro-to-machine-learning}
\end{itemize}

\subsection*{When something will not click}

\begin{itemize}[leftmargin=*]
  \item \textbf{StatQuest} (Josh Starmer) --- short, clear videos on almost every
    topic in this syllabus. Start with \emph{A Gentle Introduction to Machine
    Learning}. \textbf{[free]} \url{https://statquest.org/video-index/}
  \item \textbf{MLU-Explain} --- interactive visual explanations of train/test
    splits, bias--variance and cross-validation. \textbf{[free]}
    \url{https://mlu-explain.github.io/}
  \item \textbf{scikit-learn, Common Pitfalls} --- the official treatment of
    leakage and \texttt{random\_state}. Read it once, properly. \textbf{[free]}
    \url{https://scikit-learn.org/stable/common_pitfalls.html}
\end{itemize}

\subsection*{Course materials}

Slides, notes, lab briefs and starter notebooks are published as they are
delivered at\\
\url{https://github.com/tali7c/Applied-Machine-Learning}

Each lecture ships with \texttt{slides.pdf} and a much fuller \texttt{notes.pdf}
containing worked examples, runnable code with its real output, practice
questions and full solutions. \textbf{The notes are written so you can learn the
material without having been in the lecture} --- use them.

\vfill
\begin{center}
\small Questions about anything in this handbook: ask in class, in the lab, or
on the LMS.
\end{center}

\end{document}
